\clearpage
\hypertarget{native-483-victoria}{}
\begin{center}
\rule{0.55\linewidth}{0.4pt}\\[2pt]
Astronomische Nachrichten. Band XXXI. Nr. 740. S. 305. 22. Nov. 1850.\\[2pt]
\rule{0.55\linewidth}{0.4pt}
\end{center}
\begin{flushright}Göttingen, 31. Oktober 1850.\end{flushright}

Ich schicke Ihnen, mein teuerster Freund, meine drei Meridianbeobachtungen der \emph{Victoria}. Abgesehen davon, dass die Beobachtung in einem auch nur schwach beleuchteten Feld schwierig ist, leidet ihre Genauigkeit etwas darunter, dass die Reduktionselemente jeweils erst aus weiter entfernten Beobachtungen gewonnen werden konnten.

\begin{center}
\setlength{\tabcolsep}{3pt}
\begin{tabular}{l|r|r|r}
1850 & Göttinger m. Z. & Rektaszension & Deklination \\\hline
Sept. 30 & $10^{\rm u}\ 54^{\rm m}\ 19^{\rm s}7$ & $352^\circ\ 52'\ 43''4$ & $11^\circ\ 16'\ 48''3$ N. \\
Okt. 12 & $10\ 0\ 49.7$ & $351\ 17\ 39.6$ & $9\ 5\ 23.0$ \\
\hfill 20 & $9\ 27\ 25.9$ & $350\ 48\ 27.6$ & $7\ 45\ 38.6$ \\
\end{tabular}
\end{center}
